% =====================================================================
%  ARRIVAL Protocol — Companion Paper: AutoGen (AG2) Validation
%  DOI: 10.5281/zenodo.18893515
%  License: CC BY-NC 4.0 (text), AGPL-3.0-or-later (code)
%  Author: Mefodiy Kelevra — ORCID 0009-0003-4153-392X
% =====================================================================
\documentclass[11pt,a4paper]{article}

% ── Geometry ──────────────────────────────────────────────────────────
\usepackage[margin=2.4cm,top=2.8cm,bottom=2.8cm]{geometry}

% ── Fonts & Encoding ─────────────────────────────────────────────────
\usepackage[T1]{fontenc}
\usepackage[utf8]{inputenc}
\usepackage{lmodern}
\usepackage{microtype}

% ── Math ─────────────────────────────────────────────────────────────
\usepackage{amsmath,amssymb,amsthm}

% ── Tables & Figures ─────────────────────────────────────────────────
\usepackage{booktabs}
\usepackage{tabularx}
\usepackage{multirow}
\usepackage{graphicx}
\usepackage{float}

% ── Colours & Boxes ──────────────────────────────────────────────────
\usepackage[dvipsnames]{xcolor}
\usepackage[most]{tcolorbox}

% ── Hyperlinks ───────────────────────────────────────────────────────
\usepackage[colorlinks=true,linkcolor=NavyBlue,citecolor=OliveGreen,
            urlcolor=BrickRed]{hyperref}
\usepackage{url}

% ── Bibliography ─────────────────────────────────────────────────────
\usepackage[numbers,sort&compress]{natbib}

% ── Section formatting ───────────────────────────────────────────────
\usepackage{titlesec}
\titleformat{\section}{\Large\bfseries}{\thesection.}{0.5em}{}
\titleformat{\subsection}{\large\bfseries}{\thesubsection}{0.5em}{}
\titleformat{\subsubsection}{\normalsize\bfseries}{\thesubsubsection}{0.5em}{}

% ── Theorem-like environments (tcolorbox style) ─────────────────────
\newtcolorbox{resultbox}[1][]{
  colback=red!3, colframe=BrickRed!60!black,
  fonttitle=\bfseries, title={#1},
  boxrule=0.6pt, arc=2pt, left=6pt, right=6pt, top=4pt, bottom=4pt
}
\newtcolorbox{definitionbox}[1][]{
  colback=green!3, colframe=OliveGreen!70!black,
  fonttitle=\bfseries, title={#1},
  boxrule=0.6pt, arc=2pt, left=6pt, right=6pt, top=4pt, bottom=4pt
}

% ── Misc ─────────────────────────────────────────────────────────────
\usepackage{enumitem}
\setlist{nosep,leftmargin=*}
\usepackage{xspace}
\newcommand{\arrival}{\textsc{Arrival}\xspace}
\newcommand{\care}{\textsc{Care}\xspace}
\newcommand{\pp}{\,pp\xspace}

% ── Header / Footer ──────────────────────────────────────────────────
\usepackage{fancyhdr}
\pagestyle{fancy}
\fancyhf{}
\fancyhead[L]{\small\textit{ARRIVAL Protocol — AutoGen Companion Paper}}
\fancyhead[R]{\small\textit{DOI: 10.5281/zenodo.18893515}}
\fancyfoot[C]{\thepage}
\renewcommand{\headrulewidth}{0.4pt}

% =====================================================================
\begin{document}

% ── Title Page ───────────────────────────────────────────────────────
\begin{titlepage}
\centering
\vspace*{2cm}

{\Huge\bfseries Framework-Agnostic Validation\\
of the ARRIVAL Protocol}\\[0.8cm]
{\Large Cross-Architecture AI Coordination\\
on AutoGen (AG2)}\\[2.5cm]

{\large Mefodiy Kelevra}\\[0.3cm]
{\small ORCID:
  \href{https://orcid.org/0009-0003-4153-392X}{0009-0003-4153-392X}}\\[0.3cm]
{\small Contact: emkelvra@gmail.com}\\[2cm]

{\normalsize February 22, 2026 \enspace|\enspace Version 1.0}\\[0.5cm]
{\small DOI: \href{https://doi.org/10.5281/zenodo.18893515}%
  {10.5281/zenodo.18893515}}\\[2cm]

\rule{0.6\textwidth}{0.4pt}\\[0.8cm]

{\footnotesize
  \textbf{License} — Text: CC BY-NC 4.0 \enspace|\enspace
  Code: AGPL-3.0-or-later\\[0.4cm]
  Companion to: Kelevra, M. (2026). \textit{\arrival Protocol:
  Cross-Architecture AI Coordination through Structured Semantic Atoms.}\\[0.2cm]
  Repository:
  \url{https://github.com/mifkilla/Arrival-Protocol}
}

\vfill
\end{titlepage}

% ── Table of Contents ────────────────────────────────────────────────
\tableofcontents
\newpage

% =====================================================================
\section{Abstract}
% =====================================================================

Multi-agent AI systems increasingly rely on coordination protocols to
enable productive interaction among heterogeneous language models.  The
\arrival Protocol uses structured semantic markers (``@-atoms'') from the
DEUS.PROTOCOL~v0.5 vocabulary to achieve cross-architecture coordination
without shared training or fine-tuning.  Prior work validated \arrival
through 459 experiments on a custom Python runner.  A critical open
question is whether the protocol's coordination quality is an artifact
of the custom implementation or a genuine property of the semantic
protocol itself.

This companion paper presents 16 experiments testing \arrival on the AG2
(AutoGen) multi-agent framework, addressing the framework-agnostic
hypothesis directly.  Three experiments were conducted:
\begin{enumerate}
  \item \textbf{Baseline MCQ consensus} (5 questions, 3 agents):
    100\% accuracy, \care Resolve $= 1.000$, Meaning Debt $= 0.000$.
  \item \textbf{Adversarial robustness} under Trojan Atoms attack
    (3 questions, 4 agents including 1 saboteur): 0\% \care degradation,
    0\% false consensus, complete rejection of all 12 malicious atom
    attempts.
  \item \textbf{Head-to-head comparison} of AG2 versus native runner on
    identical inputs: 100\% answer match (5/5), \care difference $= 0.000$.
\end{enumerate}

Total experimental cost: \$0.024.  These results confirm that \arrival's
coordination quality emerges from protocol-level semantic design rather
than implementation infrastructure, establishing the first
cross-framework validation methodology for multi-agent coordination
protocols.

\textbf{Keywords:} multi-agent systems, LLM coordination,
framework-agnostic protocols, adversarial robustness, semantic atoms,
AutoGen, cross-architecture AI

% =====================================================================
\section{Introduction}\label{sec:intro}
% =====================================================================

\subsection{Motivation}

The \arrival Protocol (Kelevra, 2025--2026) demonstrates that
heterogeneous AI agents --- built on different architectures, by
different companies, with different training data --- can achieve genuine
semantic coordination through structured @-atom protocols defined by the
DEUS.PROTOCOL~v0.5 specification.  The parent paper validates this claim
through 459 experiments across 12 experimental phases, using a custom
Python runner with direct API calls to OpenRouter.

However, a natural objection arises: \textbf{does the coordination
quality depend on the specific implementation}, or is it a genuine
property of the semantic protocol?  If the custom runner's message
routing, prompt formatting, or conversation management contributes
materially to coordination outcomes, then \arrival's claims of
protocol-level coordination are weakened.  Conversely, if identical
coordination quality is achieved on an entirely different framework,
the protocol's generality is strongly supported.

This paper addresses this question by reimplementing \arrival on the AG2
(AutoGen) framework~\cite{wu2023autogen} and conducting controlled
experiments that isolate protocol effects from infrastructure effects.

\subsection{Research Questions}

\textbf{RQ1 (Framework-Agnosticism):} Does the \arrival Protocol achieve
equivalent coordination quality when executed through AG2 GroupChat
compared to a custom Python runner?

\textbf{RQ2 (Portable Robustness):} Is adversarial robustness against
Trojan Atoms attacks preserved across framework boundaries, confirming
that robustness is a protocol-level property?

\textbf{RQ3 (Methodology):} Can a reproducible methodology be
established for cross-framework validation of multi-agent coordination
protocols?

\subsection{Contributions}

\begin{enumerate}
  \item First demonstration of DEUS.PROTOCOL @-atoms functioning within
    AG2 GroupChat architecture.
  \item Empirical confirmation of the framework-agnostic hypothesis:
    identical outcomes on AG2 and native runner.
  \item Evidence that adversarial robustness is protocol-level, not
    infrastructure-level.
  \item A reusable cross-framework comparison methodology for
    multi-agent protocol validation.
\end{enumerate}

% =====================================================================
\section{Related Work}\label{sec:related}
% =====================================================================

\subsection{Multi-Agent LLM Frameworks}

\textbf{AG2/AutoGen}~\cite{wu2023autogen} is an open-source multi-agent
conversation framework originally developed at Microsoft Research.  It
provides GroupChat orchestration, configurable speaker selection, and
OpenAI-compatible API integration.

\textbf{CrewAI} and \textbf{LangGraph} offer alternative multi-agent
orchestration paradigms.  CrewAI emphasizes role-based agent
specialization; LangGraph emphasizes graph-based workflow definition.
Testing \arrival on these frameworks remains future work.

\subsection{LLM Collective Intelligence}

Riedl et al.\ (2025, arXiv:2510.05174) study LLM collective
intelligence through deliberation.  Ashery et al.\ (2025, Science
Advances) demonstrate that LLM group accuracy exceeds individual
accuracy under certain conditions.  Du et al.~\cite{du2024debate} show
that multi-agent debate improves factuality and reasoning.  \arrival
differs from debate-based approaches by imposing structured semantic
atoms rather than free-form debate.

\subsection{Protocol Portability}

The concept of protocol portability --- whether a coordination protocol
works identically across different execution environments --- is
well-established in distributed systems (e.g., Paxos, Raft).  However,
in the LLM multi-agent domain, protocol portability has not been
systematically studied.

\subsection{Adversarial Robustness in MAS}

La Malfa et al.\ (2025, arXiv:2505.21298, NeurIPS 2025) critically
examine claims about multi-agent LLM system robustness, arguing that
many reported improvements are artifacts of experimental design.  This
critique motivates our cross-framework adversarial testing.

% =====================================================================
\section{Methodology}\label{sec:method}
% =====================================================================

\subsection{Experimental Design}

All experiments share the following configuration:

\begin{table}[H]
\centering\small
\caption{Shared experimental configuration.}\label{tab:config}
\begin{tabular}{@{}ll@{}}
\toprule
\textbf{Parameter} & \textbf{Value} \\
\midrule
Framework       & AG2 (AutoGen) v0.4.x \\
API Provider    & OpenRouter \\
Models          & DeepSeek V3, Llama 3.3 70B, Qwen 2.5 72B \\
Temperature     & 0.3 (all models) \\
Protocol        & DEUS.PROTOCOL v0.5 (66 standard atoms) \\
Coordination metric  & \care Resolve \\
Semantic health metric & Meaning Debt \\
\bottomrule
\end{tabular}
\end{table}

Each agent is instantiated as an \texttt{autogen.AssistantAgent} with
model-specific \texttt{llm\_config} pointing to OpenRouter.  Agents are
grouped in an \texttt{autogen.GroupChat} with
\texttt{speaker\_selection\_method="round\_robin"} and a
\texttt{GroupChatManager} orchestrating the conversation.

\subsection{CARE Scoring}

\care Resolve is a composite metric from MEANING-CRDT
v1.1~\cite{kelevra2026crdt}:
\[
  \hat{\mathbf{v}} = \frac{\sum_{i} w_i \cdot \mathbf{v}_i}
                          {\sum_{i} w_i}
\]
where $\mathbf{v}_i$ is the position vector of agent $i$ and $w_i$ is
its epistemic weight.  \care decomposes into four dimensions scored
0.0--1.0: \textbf{C}onsistency, \textbf{A}greement, \textbf{R}esolve,
\textbf{E}vidence.

\textbf{Meaning Debt} $\mathrm{MD}_i(T) = \sum_{k=1}^{T} L_i(k)$
accumulates semantic losses per agent per round.  $\mathrm{MD} = 0$
indicates perfect semantic health.

\subsection{Experiment 1: Baseline MCQ Consensus}

\begin{table}[H]
\centering\small
\caption{Experiment 1 parameters.}\label{tab:exp1params}
\begin{tabular}{@{}ll@{}}
\toprule
Agents       & 3 (DeepSeek V3, Llama 3.3, Qwen 2.5) \\
Rounds       & 4 per question \\
Questions    & 5 (Phase~7 hardest set) \\
API calls    & ${\sim}60$ \\
Cost         & \$0.0059 \\
\bottomrule
\end{tabular}
\end{table}

\subsection{Experiment 2: Adversarial Robustness}

\begin{table}[H]
\centering\small
\caption{Experiment 2 parameters.}\label{tab:exp2params}
\begin{tabular}{@{}ll@{}}
\toprule
Agents       & 4 (3 honest + 1 saboteur: Gemma-2 27B) \\
Strategy     & Trojan Atoms (fabricated @-atoms with corrupted semantics) \\
Questions    & 3 (with matched 3-agent honest controls) \\
Experiments  & 6 (3 adversarial + 3 control) \\
Cost         & \$0.0095 \\
\bottomrule
\end{tabular}
\end{table}

\subsection{Experiment 3: Cross-Framework Comparison}

\begin{table}[H]
\centering\small
\caption{Experiment 3: AG2 vs.\ Native Runner.}\label{tab:exp3params}
\begin{tabular}{@{}lll@{}}
\toprule
\textbf{Parameter} & \textbf{AG2} & \textbf{Native} \\
\midrule
Framework     & AG2 GroupChat          & Custom Python (direct HTTP) \\
Models        & DeepSeek V3, Llama, Qwen & Identical \\
Prompts       & DEUS.PROTOCOL v0.5     & Identical \\
Temperature   & 0.3                     & Identical \\
Questions     & 5                       & Identical \\
Cost          & \multicolumn{2}{c}{\$0.0086 (both)} \\
\bottomrule
\end{tabular}
\end{table}

\textbf{Pre-registered threshold}: Same answer rate $\geq 80\%$ AND
$|\text{CARE(AG2)} - \text{CARE(Native)}| < 0.1$.

% =====================================================================
\section{Results}\label{sec:results}
% =====================================================================

\subsection{Experiment 1: Baseline Consensus on AG2}

\begin{table}[H]
\centering\small
\caption{Experiment 1 accuracy (5 questions).}\label{tab:exp1acc}
\begin{tabular}{@{}llcccccc@{}}
\toprule
\textbf{Q} & \textbf{Domain} & \textbf{Correct} &
  \textbf{DeepSeek} & \textbf{Llama} & \textbf{Qwen} &
  \textbf{MV} & \textbf{\arrival} \\
\midrule
SCI\_07  & Science  & B & \checkmark & \checkmark & \checkmark & \checkmark & \checkmark \\
HIS\_08  & History  & B & \checkmark & \checkmark & \checkmark & \checkmark & \checkmark \\
LOG\_09  & Logic    & C & $\times$   & $\times$   & \checkmark & \checkmark & \checkmark \\
LAW\_06  & Law      & B & \checkmark & \checkmark & \checkmark & \checkmark & \checkmark \\
TECH\_10 & Tech     & B & \checkmark & \checkmark & \checkmark & \checkmark & \checkmark \\
\bottomrule
\end{tabular}
\end{table}

\textbf{Solo accuracy}: DeepSeek V3: 4/5 (80\%), Llama 3.3: 4/5 (80\%),
Qwen 2.5: 5/5 (100\%).
\textbf{MV}: 5/5 (100\%).
\textbf{\arrival on AG2}: 5/5 (100\%).

\begin{table}[H]
\centering\small
\caption{\care metrics for Experiment~1.}\label{tab:exp1care}
\begin{tabular}{@{}lrrr@{}}
\toprule
\textbf{Question} & \textbf{\care Resolve} & \textbf{Meaning Debt} &
  \textbf{Unique Atoms} \\
\midrule
SCI\_07   & 1.000 & 0.000 & 6 \\
HIS\_08   & 1.000 & 0.000 & 7 \\
LOG\_09   & 1.000 & 0.000 & 6 \\
LAW\_06   & 1.000 & 0.000 & 7 \\
TECH\_10  & 1.000 & 0.000 & 6 \\
\midrule
\textbf{Mean} & \textbf{1.000} & \textbf{0.000} & \textbf{6.4} \\
\bottomrule
\end{tabular}
\end{table}

All questions achieved perfect \care Resolve and zero Meaning Debt.

\subsection{Experiment 2: Adversarial Robustness on AG2}

\begin{table}[H]
\centering\small
\caption{\care degradation under Trojan Atoms attack.}\label{tab:exp2care}
\begin{tabular}{@{}lrrlll@{}}
\toprule
\textbf{Q} & \textbf{\care (Ctrl)} & \textbf{\care (Adv)} &
  \textbf{Degrad.} & \textbf{Correct?} & \textbf{False Cons.} \\
\midrule
SCI\_07 & 1.000 & 1.000 & 0.000 & Yes/Yes & No \\
HIS\_08 & 1.000 & 1.000 & 0.000 & Yes/Yes & No \\
LOG\_09 & 1.000 & 1.000 & 0.000 & Yes/Yes & No \\
\midrule
\textbf{Mean} & \textbf{1.000} & \textbf{1.000} & \textbf{0.000} &
  \textbf{3/3} & \textbf{0/3} \\
\bottomrule
\end{tabular}
\end{table}

\begin{table}[H]
\centering\small
\caption{Trojan Atom rejection rates.}\label{tab:trojan}
\begin{tabular}{@{}lrrr@{}}
\toprule
\textbf{Trojan Strategy} & \textbf{Used} & \textbf{Adopted} &
  \textbf{Rate} \\
\midrule
Fabricated \texttt{@EVIDENCE}         & 3 & 0 & 0.0\% \\
False \texttt{@CONFIDENCE(0.95)}      & 3 & 0 & 0.0\% \\
Misrepresented \texttt{@SYNTHESIS}    & 3 & 0 & 0.0\% \\
False \texttt{@CONSENSUS}             & 3 & 0 & 0.0\% \\
\midrule
\textbf{Total} & \textbf{12} & \textbf{0} & \textbf{0.0\%} \\
\bottomrule
\end{tabular}
\end{table}

The saboteur (Gemma-2 27B) introduced novel malicious atoms
(\texttt{@AUTHORITY\_OVERRIDE}, \texttt{@CONSENSUS\_BYPASS}) not present
in the DEUS.PROTOCOL vocabulary.  Honest agents ignored these entirely.
Recovery rate: 3/3 (100\%).

\subsection{Experiment 3: AG2 vs.\ Native Runner}

\begin{table}[H]
\centering\small
\caption{Answer comparison: AG2 vs.\ Native.}\label{tab:exp3answers}
\begin{tabular}{@{}lllcll@{}}
\toprule
\textbf{Q} & \textbf{AG2} & \textbf{Native} & \textbf{Match} &
  \textbf{AG2 Corr.} & \textbf{Nat.\ Corr.} \\
\midrule
SCI\_07  & B & B & Yes & Yes & Yes \\
HIS\_08  & B & B & Yes & Yes & Yes \\
LOG\_09  & C & C & Yes & Yes & Yes \\
LAW\_06  & B & B & Yes & Yes & Yes \\
TECH\_10 & B & B & Yes & Yes & Yes \\
\bottomrule
\end{tabular}
\end{table}

\begin{table}[H]
\centering\small
\caption{Atom usage comparison.}\label{tab:exp3atoms}
\begin{tabular}{@{}lrrr@{}}
\toprule
\textbf{Metric} & \textbf{AG2} & \textbf{Native} & \textbf{$\Delta$} \\
\midrule
Mean atom count/dialogue & 8.6   & 6.4   & +2.2 \\
Mean emergent atoms      & 1.8   & 0.2   & +1.6 \\
\care Resolve            & 1.000 & 1.000 & 0.000 \\
Accuracy                 & 5/5   & 5/5   & 0 \\
\bottomrule
\end{tabular}
\end{table}

\begin{resultbox}[Framework-Agnostic Verdict]
Same answer rate $= 100\%$ ($\geq 80\%$ threshold met).\\
\care difference $= 0.000$ ($< 0.1$ threshold met).\\
\textbf{Framework-agnostic hypothesis CONFIRMED.}
\end{resultbox}

\subsection{Aggregate Summary}

\begin{table}[H]
\centering\small
\caption{Aggregate results across all experiments.}\label{tab:aggregate}
\begin{tabular}{@{}lrrrr@{}}
\toprule
\textbf{Metric} & \textbf{Exp.~1} & \textbf{Exp.~2} &
  \textbf{Exp.~3 AG2} & \textbf{Exp.~3 Nat.} \\
\midrule
Accuracy      & 5/5 (100\%) & 3/3 (100\%) & 5/5 (100\%) & 5/5 (100\%) \\
\care Resolve & 1.000       & 1.000       & 1.000       & 1.000 \\
Meaning Debt  & 0.000       & 0.000       & 0.000       & 0.000 \\
Atoms (mean)  & 6.4         & ---         & 8.6         & 6.4 \\
Cost (USD)    & 0.0059      & 0.0095      & 0.0086      & (incl.) \\
\bottomrule
\end{tabular}
\end{table}

\textbf{Grand total}: 16 experiments, \$0.024 total cost.

% =====================================================================
\section{Discussion}\label{sec:discussion}
% =====================================================================

\subsection{Framework-Agnosticism Confirmed}

The central finding is that the \arrival Protocol produces identical
coordination outcomes on AG2 and on a custom Python runner: 100\%
answer match, identical \care Resolve (1.000), identical Meaning Debt
(0.000).  The only difference --- AG2 producing slightly more @-atoms
per dialogue (8.6 vs.\ 6.4) --- is a surface-level formatting effect
that does not impact coordination quality.

\subsection{Protocol-Level Adversarial Robustness}

Experiment~2 demonstrates that \arrival's resistance to Trojan Atoms
attacks is preserved on AG2.  The saboteur deployed 12 Trojan Atom
instances achieving 0\% adoption rate and 0\% \care degradation.  This
partially addresses the critique of La Malfa et al.\ (2025): by
demonstrating identical robustness on two independent frameworks, we
provide evidence that \arrival's adversarial resistance is a genuine
property of the protocol.

\subsection{Methodological Contribution}

The cross-framework comparison methodology established here provides a
reusable template:
\begin{enumerate}
  \item Fix all variables except the framework
  \item Run identical experiments on both frameworks
  \item Compare using pre-registered thresholds
  \item Report both outcome and process metrics
\end{enumerate}

\subsection{Cost Efficiency}

The 16 experiments cost \$0.024 total, averaging \$0.0015 per
experiment --- consistent with the parent paper's cost efficiency.

% =====================================================================
\section{Limitations}\label{sec:limitations}
% =====================================================================

\begin{enumerate}
  \item \textbf{Sample size}: 5 questions for Experiments 1 and 3,
    3 questions for Experiment 2.
  \item \textbf{Single alternative framework}: Only AG2 tested.
  \item \textbf{Model non-determinism}: Even at temperature 0.3,
    outputs are stochastic.
  \item \textbf{MCQ-only}: Framework-agnostic hypothesis not tested
    on open-ended tasks.
  \item \textbf{English-only}: Cross-linguistic performance not tested.
  \item \textbf{Ceiling effects}: Perfect \care scores may mask real
    framework differences on harder tasks.
\end{enumerate}

% =====================================================================
\section{Conclusion}\label{sec:conclusion}
% =====================================================================

This companion paper presents 16 experiments testing the \arrival
Protocol on AG2 (AutoGen).  Three findings emerge:

\begin{enumerate}
  \item \textbf{\arrival works on AG2.} 100\% accuracy with
    \care $= 1.000$ and Meaning Debt $= 0.000$ on AG2 GroupChat.
  \item \textbf{Adversarial robustness is protocol-level.} Trojan Atoms
    attack achieves 0\% success rate on AG2, with all 12 malicious atom
    instances rejected.
  \item \textbf{Framework-agnosticism confirmed.} AG2 and native runner
    produce identical answers (5/5 match) with identical \care scores
    ($\Delta = 0.000$).
\end{enumerate}

These results have practical implications: organizations can deploy
\arrival on their existing multi-agent infrastructure without
coordination quality loss.

% =====================================================================
% Bibliography
% =====================================================================

\begin{thebibliography}{11}

\bibitem{kelevra2026arrival}
Kelevra, M. (2025--2026).
\newblock ARRIVAL Protocol: Cross-Architecture AI Coordination through
  Structured Semantic Atoms.
\newblock \textit{Preprint.}
\newblock ORCID: 0009-0003-4153-392X.

\bibitem{kelevra2026crdt}
Kelevra, M. (2026).
\newblock MEANING-CRDT v1.1: A Mathematical Framework for Semantic
  Convergence in Multi-Agent Systems.
\newblock DOI: \href{https://doi.org/10.5281/zenodo.18702383}%
  {10.5281/zenodo.18702383}.

\bibitem{wu2023autogen}
Wu, Q., Bansal, G., Zhang, J., et al. (2023).
\newblock AutoGen: Enabling Next-Gen LLM Applications via Multi-Agent
  Conversation.
\newblock \textit{arXiv:2308.08155.}

\bibitem{du2024debate}
Du, Y., Li, S., Torralba, A., Tenenbaum, J.~B., \& Mordatch, I. (2024).
\newblock Improving Factuality and Reasoning in Language Models through
  Multiagent Debate.
\newblock \textit{ICML 2024.}

\bibitem{riedl2025collective}
Riedl, C., et al. (2025).
\newblock LLM Collective Intelligence through Deliberation.
\newblock \textit{arXiv:2510.05174.}

\bibitem{ashery2025social}
Ashery, S., et al. (2025).
\newblock LLM Group Accuracy Exceeds Individual Accuracy.
\newblock \textit{Science Advances.}

\bibitem{chan2024chateval}
Chan, C.-M., et al. (2024).
\newblock ChatEval: Towards Better LLM-based Evaluators through
  Multi-Agent Debate.
\newblock \textit{arXiv:2308.07201.}

\bibitem{lamalfa2025limits}
La~Malfa, E., et al. (2025).
\newblock A Critical Examination of Multi-Agent LLM System Claims.
\newblock \textit{arXiv:2505.21298.} NeurIPS 2025.

\bibitem{vlmtoc2026}
VLM-TOC. (2026).
\newblock Theory of Collaboration for Vision-Language Models.
\newblock \textit{arXiv:2601.20641.}

\end{thebibliography}

% =====================================================================
\appendix
\section{Reproducibility}

\begin{itemize}
  \item \textbf{Python:} 3.10+
  \item \textbf{AG2:} \texttt{pip install -U "ag2[openai]"}
  \item \textbf{API:} OpenRouter account with API key
  \item \textbf{Models:} \texttt{deepseek/deepseek-chat-v3-0324},
    \texttt{meta-llama/llama-3.3-70b-instruct},
    \texttt{qwen/qwen-2.5-72b-instruct},
    \texttt{google/gemma-2-27b-it}
\end{itemize}

\section{DEUS.PROTOCOL v0.5 Atom Subset}

\begin{table}[H]
\centering\small
\caption{Atoms used in AG2 experiments.}\label{tab:atoms}
\begin{tabular}{@{}lll@{}}
\toprule
\textbf{Atom} & \textbf{Category} & \textbf{Round} \\
\midrule
\texttt{@ANSWER}        & Response   & 1, 4 \\
\texttt{@CONFIDENCE(x)} & Meta       & All \\
\texttt{@EVIDENCE}       & Reasoning  & 1, 3 \\
\texttt{@CHALLENGE}      & Critique   & 2 \\
\texttt{@AGREE}          & Alignment  & 2, 3 \\
\texttt{@DISAGREE}       & Alignment  & 2 \\
\texttt{@UPDATE}         & Revision   & 3 \\
\texttt{@MAINTAIN}       & Revision   & 3 \\
\texttt{@STRENGTHEN}     & Revision   & 3 \\
\texttt{@SYNTHESIS}      & Integration & 4 \\
\texttt{@CONSENSUS}      & Resolution & 4 \\
\texttt{@RESOLVE}        & Resolution & 4 \\
\bottomrule
\end{tabular}
\end{table}

\end{document}
